\title{Direct Preference Optimization: Your Language Model is Secretly a Reward Model}

\begin{abstract}
Large-scale unsupervised language models (LMs) acquire extensive general knowledge and demonstrate certain reasoning capabilities; however, their behaviors are challenging to direct due to the inherent lack of supervision in their training process. Common approaches to address this challenge involve gathering human judgments on the comparative quality of outputs produced by the model and subsequently fine-tuning the unsupervised LM to conform to these human preferences, frequently utilizing reinforcement learning from human feedback (RLHF). Despite their effectiveness, RLHF-based techniques are typically complicated and prone to instability, as they involve a two-step process: first, constructing a reward model to capture human preferences, and second, using reinforcement learning to adjust the large LM so that it optimizes for the predicted reward without significantly diverging from its initial parameters. In this work, we propose a novel reward model parameterization within RLHF that permits the derivation of the optimal policy in closed form. This reformulation allows us to address the conventional RLHF objective by minimizing a simple classification loss, leading to an algorithm we term \textit{Direct Preference Optimization} (DPO). DPO offers improved stability and efficiency, removing the necessity for sampling from the LM during fine-tuning and obviating extensive hyperparameter searches. Through empirical studies, we demonstrate that DPO is capable of adapting LMs to human preferences on par with or surpassing current alternatives. Specifically, DPO provides superior sentiment control over PPO-based RLHF and delivers equal or better performance on tasks like summarization and single-turn dialogue, all while being easier to deploy and requiring fewer computational resources.
\end{abstract}

\section{Introduction}
Unsupervised language models (LMs) that are trained on massive corpora have demonstrated a range of unexpected abilities~\citep{chowdhery2022palm, brown2020language, touvron2023llama,bubeck2023sparks}. Despite these capabilities, such models are built on data reflecting human outputs with diverse objectives, preferences, and proficiencies. Inevitably, certain intentions and competencies captured in the data may not be optimal for emulation; for instance, although it is advantageous for an AI code assistant to \textit{recognize} frequent programming errors to enable effective correction, we ultimately prefer model generations to reflect the (sometimes uncommon) superior coding proficiency included in the training set. Likewise, a language model’s \textit{awareness} of a widespread misconception held by half the population should not entail endorsing that false belief in half of all pertinent queries. This highlights the essential distinction between a model’s vast \textit{knowledge and competencies} and its target \emph{behaviors and outputs}: configuring the latter is pivotal for constructing AI systems that are reliable, effective, and amenable to user control \citep{ouyang2022training}. Standard methodologies frequently guide LMs towards alignment with human values through reinforcement learning (RL), but, as we demonstrate, the RL-derived optimization objectives used by current techniques can, in fact, be exactly matched using a simple binary cross-entropy loss, leading to a substantial streamlining of preference learning processes.

Broadly speaking, prevailing approaches incorporate human-aligned behaviors into language models via carefully selected preference datasets which capture the kinds of actions perceived as beneficial and safe. This phase of learning from preferences follows a preliminary unsupervised pre-training stage conducted on large-scale textual data. Although direct supervised fine-tuning using exemplary human-generated outputs offers a straightforward avenue for preference learning, the most effective techniques belong to the reinforcement learning from human or AI feedback family (RLHF/RLAIF; \citep{christiano2017deep,bai2022constitutional}). In RLHF, a reward model is trained to interpret a database of human preferences, and the LM is subsequently optimized with RL so that its policy favors high-reward outputs without straying unduly from the original distribution. While RLHF models achieve exceptional proficiency in dialogue and coding applications, their training pipelines are considerably more involved than those for supervised learning—necessitating the training of several LMs and integrating sampling from the LM policy during optimization—which results in significant additional computational burden.

This work presents a method for optimizing language models according to human preferences without relying on explicit reward modeling or reinforcement learning. We introduce \textit{{Direct Preference Optimization} (DPO)}, a technique that implicitly targets the same objective as traditional RLHF algorithms—specifically, reward maximization under a KL-divergence regularization—but with a more accessible and easily trainable framework. Conceptually, DPO updates work by raising the relative log-likelihood of favored responses over those not preferred; crucially, each update is weighted on a per-sample basis to prevent the kind of model collapse we observe with naive probability-ratio objectives. DPO, like prior approaches, draws upon a theoretical preference model (e.g., the Bradley-Terry model; \cite{bradley1952rankanalysis}) for quantifying the concordance between a candidate reward function and observed preference data. Distinctly, whereas prior RLHF methods utilize the preference model to derive a loss for training a reward model and subsequently optimize a policy against this learned reward, DPO leverages a variable transformation, directly expressing the preference loss in terms of the policy. Given a dataset of human-labeled preferences among generated responses, DPO enables policy training through a simple binary cross-entropy loss, effectively discovering the optimal policy for an implicit reward shaped by the preference data.

Our principal contribution is the development of Direct Preference Optimization (DPO), an RL-free, streamlined approach for preference-based language model training. Empirical evaluations indicate that DPO matches or exceeds the performance of existing methodologies, such as PPO-based RLHF, in applications including sentiment transfer, text summarization, and conversation, even when scaling up to language models containing as many as 6 billion parameters.

\section{Related Work}

Self-supervised language models that are scaled up can address certain tasks without explicit training data (zero-shot) \citep{radford2019language} or when provided with a handful of example prompts (few-shot) \citep{gpt3,megatron,chowdhery2022palm}. Yet, their capability on specific downstream applications and their adherence to user intentions are markedly enhanced when they are further trained—fine-tuned—on corpora consisting of instructions and completions authored by humans \citep{mishra-etal-2022-cross,sanh2022multitask,chung2022scaling,thoppilan2022lamda}. This process, referred to as `instruction-tuning,’ facilitates broader generalization by large language models (LLMs) to handle instructions not seen during tuning, thus improving their general utility \citep{chung2022scaling}. Even with its demonstrated success, collecting \textit{relative} evaluations of response quality from humans is typically more practical than gathering expert-crafted examples. As a result, later efforts have refined LLMs with datasets reflecting human preferences, which has led to performance gains across tasks such as translation \citep{kreutzer-etal-2018-reliability}, summarization \citep{stiennon2022learning,ziegler2020finetuning}, narrative generation \citep{ziegler2020finetuning}, and compliance with user instructions \citep{ouyang2022training,ramamurthy2023is}. These techniques initially learn a neural reward function that aligns with preference datasets under models like the Bradley-Terry preference framework \citep{bradley1952rankanalysis}. Subsequently, the language model is fine-tuned to maximize these learned rewards using reinforcement learning approaches, most notably REINFORCE \citep{williams1992reinforce}, proximal policy optimization (PPO; \cite{schulman2017proximal}), or related variants \citep{ramamurthy2023is}. In parallel, another research stream has made use of instruction-following LLMs, refined via human feedback, to synthesize new preference data targeting characteristics like safety or harmlessness \citep{bai2022constitutional}, often relying solely on minimal human oversight conveyed through textual rubrics to guide model annotations. These developments integrate methodologies from two domains: one dedicated to applying reinforcement learning techniques to language model training for diverse objectives~\citep{Ranzato2015SequenceLT,paulus2018a,wu2018learning}, and another centered on learning from human preference signals \citep{christiano2017deep,kupcsik2018learning}. While relative human preferences offer practical advantages, reinforcement learning-based fine-tuning for large language models presents considerable implementation challenges; in response, this study introduces a theoretically principled alternative for optimizing relative preferences that circumvents the need for reinforcement learning.

Beyond language applications, the task of learning policies from preferences has received significant attention in both bandit and reinforcement learning paradigms, leading to a variety of methodological developments. When contextual bandits are trained using action preferences or rankings in lieu of explicit rewards, the problem formulation is termed a contextual dueling bandit (CDB; \cite{yue2012karmed,dudik2015contextual}). Without access to concrete rewards, theoretical analyses of CDBs typically replace the idea of an optimal policy with that of a \textit{von Neumann winner}, defined as a policy that is expected to achieve at least a 50\% success rate in pairwise comparisons against any alternative policy \citep{dudik2015contextual}. Importantly, CDB frameworks obtain preference labels in an online manner, whereas human preference learning most often relies on a finite, static dataset consisting of offline, preference-labelled action pairs \citep{yan2022human}. Likewise, in \textit{preference-based RL} (PbRL), agents are guided by binary feedback derived from some unknown 'scoring' mechanism, as opposed to receiving reward signals directly \citep{BusaFekete2014,ruiz2023dueling}. There exists a spectrum of PbRL algorithms—some equipped to leverage off-policy preference data—but these approaches commonly involve first reconstructing a hidden scoring or reward model, followed by policy optimization against it \citep{jain2013learning,BusaFekete2014,christiano2017deep,sadigh2017active,kupcsik2018learning}. In contrast, we introduce a streamlined, one-step policy learning procedure that directly trains a policy to align with stated preferences, bypassing the need for explicit reward estimation.

\section{Preliminaries}\label{section:prelims}

We briefly outline the RLHF methodology as formulated by \citeauthor{ziegler2020finetuning}, subsequently extended in works such as \citep{stiennon2022learning, bai2022training, ouyang2022training}. This process is conventionally divided into three distinct stages: (1) supervised fine-tuning (SFT), (2) preference collection and reward model learning, and (3) reinforcement learning-based optimization.

\textbf{SFT}: The RLHF process typically commences by adapting a pre-trained language model to downstream objectives (such as dialogue or summarization) through supervised fine-tuning on curated datasets, resulting in an initial model denoted as $\pi^\text{SFT}$.

\textbf{Reward Modelling Phase}: During this second phase, the fine-tuned model $\pi^\text{SFT}$ is queried with input prompts $x$, generating answer pairs $(y_1, y_2)\sim \pi^\text{SFT}(y \mid x)$. Human annotators are then presented with these pairs and asked to indicate which answer they prefer, denoted $y_w\succ y_l \mid x$, where $y_w$ and $y_l$ correspond to the chosen and non-chosen answers from the pair, respectively. Although preferences are presumed to be governed by an underlying reward function $r^*(y, x)$ that remains unobserved, various modeling approaches have been proposed for this setting. The Bradley-Terry (BT) model \cite{bradley1952rankanalysis} is frequently employed (with generalizations to Plackett-Luce models \citep{plackett1975analysis, luce2012individual} feasible when multi-way rankings are available). Under the BT framework, the distribution of human preferences $p^*$ is specified as follows:

\begin{equation}\label{eq:bradley-terry}
    p^*(y_1\succ y_2 \mid x)=\frac{\exp\left(r^*(x, y_1)\right)}{\exp\left(r^*(x, y_1)\right) + \exp\left(r^*(x, y_2)\right)}.
\end{equation}

Given a fixed collection of preference data $\mathcal{D}=\bigl\{x^{(i)}, y_w^{(i)}, y_l^{(i)}\bigr\}_{i=1}^N$, sampled from $p^*$, a reward function $r_{\phi}(x, y)$ can be parameterized and its parameters can be determined using maximum likelihood estimation. By casting this task as a binary classification problem, the associated loss function corresponds to the negative log-likelihood:

\begin{equation}\label{eq:reward_model}
    \mathcal{L}_R(r_{\phi}, \mathcal{D}) = -\mathbb{E}_{(x, y_w, y_l)\sim \mathcal{D}}\bigl[\log \sigma(r_{\phi}(x, y_w)- r_{\phi}(x, y_l))\bigr]
\end{equation}

where $\sigma$ denotes the logistic sigmoid function. For large language models, $r_{\phi}(x, y)$ is commonly initialized from the SFT policy $\pi^\text{SFT}(y \mid x)$, typically by appending a linear transformation atop the output of the final transformer layer, yielding a single scalar reward prediction per input \cite{ziegler2020finetuning}. To reduce reward variance, previous studies have normalized the reward values so that $\mathbb{E}_{x,y\sim \mathcal{D}}\left[r_\phi(x, y)\right] = 0$ across all $x$.

\textbf{RL Fine-Tuning Phase}: Within the reinforcement learning phase, the trained reward function delivers evaluative feedback to the language model. Building on earlier methods~\citep{jaques2017sequence, jaques2020human}, the fine-tuning objective is formulated as

\begin{equation}\label{eq:RL}
\max_{\pi_{\theta}}  \mathbb{E}_{x\sim \mathcal{D}, y\sim \pi_{\theta}(y \mid x)}\bigl[r_{\phi}(x, y)\bigr] - \beta\mathbb{D}_{\textrm{KL}}\bigl[\pi_{\theta}(y\mid x)\mid \mid \pi_\text{ref}(y\mid x)\bigr],
\end{equation}

where $\beta$ regulates the regularization penalty between the learned policy $\pi_\theta$ and a reference policy $\pi_\text{ref}$—typically set to the initial SFT policy $\pi^\text{SFT}$. Generally, $\pi_\theta$ is initialized as $\pi^\text{SFT}$. This KL regularization term plays a crucial role, ensuring the policy remains within the regime where the reward model remains accurate, preserves response diversity, and mitigates mode collapse into a handful of over-optimized outputs. Since sequence generation involves discrete decisions, this objective is non-differentiable, necessitating the use of reinforcement learning algorithms. The prevalent approach \citep{ziegler2020finetuning, stiennon2022learning, bai2022training, ouyang2022training} is to define the objective function as ${r(x, y) = r_{\phi}(x, y) -\beta (\log \pi_{\theta}(y\mid x) - \log \pi_\text{ref}(y\mid x))}$, which is then optimized with PPO \cite{schulman2017proximal}. 

\section{Direct Preference Optimization}\label{sec:DPO}

Addressing the limitations of reinforcement learning methods on large-scale tasks such as language model fine-tuning, we seek a more direct strategy for policy optimization utilizing human preference data. In contrast to standard RLHF pipelines—which first learn an explicit reward model and subsequently optimize the policy with reinforcement learning—our method hinges on a specific parameterization of the reward model that facilitates direct recovery of the optimal policy in closed-form, thereby bypassing any iterative RL procedure. As elaborated in the following section, our central idea is to exploit a tractable, analytical mapping from the reward space to the space of optimal policies, allowing a reparameterization that recasts the loss function from rewards to policies. This change-of-variables technique circumvents the need to explicitly fit an independent reward model, yet remains consistent with popular formulations of human preference,

\textbf{Derivation of the DPO Objective.} The starting point is the reinforcement learning objective, Eq.~\ref{eq:RL}, considered in prior research, which is formulated with a general reward function $r$. Building on previous findings~\citep{peters2007reinforcement, peng2019advantage, korbak2022reinforcement, go2023aligning}, it can be readily shown that the solution maximizing the reward under the KL constraint, as given by Eq.~\ref{eq:RL}, is given by:

\begin{equation}\label{eq:op_policy}
    \pi_r(y\mid x) = \frac{1}{Z(x)}\pi_\text{ref}(y\mid x)\exp\left(\frac{1}{\beta}r(x, y)\right),
\end{equation}

where the normalizing constant, $Z(x) =\sum_{y}\pi_\text{ref}(y\mid x)\exp\left(\frac{1}{\beta}r(x, y)\right)$, ensures a valid probability distribution. Additional details on this derivation can be found in Appendix~\ref{app:derivation1}. Despite substituting the maximum likelihood estimate $r_{\phi}$ for the true reward function $r^*$, estimating the partition function $Z(x)$ remains computationally intensive~\citep{korbak2022reinforcement, go2023aligning}, limiting the practicality of this form. To address this, we can reformulate Eq.~\ref{eq:op_policy} to present the reward as a function of the optimal policy $\pi_r$, the reference policy $\pi_\text{ref}$, and the partition function $Z(\cdot)$. By taking logarithms of both sides in Eq.~\ref{eq:op_policy} and rearranging, we find:

\begin{equation}\label{eq:main_eq}
    r(x,y) =\beta \log \frac{\pi_r(y\mid x)}{\pi_\text{ref}(y\mid x)} + \beta \log Z(x).
\end{equation}

This transformation can be similarly applied to the ground-truth reward $r^*$ and its associated optimal policy $\pi^*$. Importantly, the Bradley-Terry model evaluates only differences in reward for two possible outputs; specifically, ${p^*(y_1 \succ y_2 \mid x) = \sigma(r^*(x, y_1) - r^*(x, y_2))}$. By substituting the expression for $r^*(x,y)$ from Eq.~\ref{eq:main_eq} into Eq.~\ref{eq:bradley-terry}, terms involving the partition function $Z(x)$ cancel out, resulting in a representation of the preference probability that is solely dependent on the optimal policy $\pi^*$ and the reference policy $\pi_\text{ref}$. Therefore, under the Bradley-Terry model, the optimal RLHF policy $\pi^*$ is governed by the following preference relation:

\begin{equation}\label{eq:objective}
    p^*(y_1\succ y_2 \mid x)=\frac{1}{1 + \exp\left(\beta \log \frac{\pi^*(y_2\mid x)}{\pi_\text{ref}(y_2\mid x)} - \beta \log \frac{\pi^*(y_1\mid x)}{\pi_\text{ref}(y_1\mid x)}\right)}
\end{equation}

For full derivation, refer to Appendix~\ref{app:derivation2}. Although Eq.~\ref{eq:objective} utilizes the Bradley-Terry model, analogous results can be established for more general models, such as the Plackett-Luce family~\citep{plackett1975analysis, luce2012individual}, with further details in Appendix~\ref{app:plackett_luce_models}.

By now expressing the likelihood of human preferences directly in terms of the optimal policy (rather than via a reward model), we are able to pose a maximum likelihood objective for a parameterized policy $\pi_\theta$. In direct analogy to traditional reward modeling (see Eq.~\ref{eq:reward_model}), the policy learning objective becomes:

\begin{equation}\label{eq:optimum_model}
    \mathcal{L}_\text{DPO}(\pi_{\theta}; \pi_\text{ref}) = -\mathbb{E}_{(x, y_w, y_l)\sim \mathcal{D}}\left[\log \sigma \left(\beta \log \frac{\pi_{\theta}(y_w\mid x

In this formulation, we approximate an implicit reward via an alternative parameterization, wherein the optimal policy coincides with $\pi_\theta$. Notably, the approach mirrors the process of fitting a reparametrized Bradley-Terry model, which confers desirable theoretical attributes, including consistency under suitable assumptions on the distribution of preference data \cite{bong2022generalized}. Additional theoretical aspects of DPO, particularly in comparison with related literature, are elaborated in Section~\ref{sec:theory}.

\textbf{What is the mechanism behind the DPO update?} To gain mechanistic insight into DPO, it is informative to examine the gradient of the objective $\mathcal{L}_\text{DPO}$. The partial derivative with respect to $\theta$ takes the form:

\begin{multline*}\label{eq:gradient}
    \nabla_\theta \mathcal{L}_\text{DPO}(\pi_\theta;\pi_\text{ref}) = \\ -\beta\mathbb{E}_{(x, y_w, y_l) \sim \mathcal{D}} \bigg[\underbrace{\sigma(\hat{r}_\theta(x, y_l) - \hat{r}_\theta (x, y_w))}_\text{higher weight when reward estimate is wrong}\bigg[\underbrace{\nabla_\theta\log \pi(y_w \mid x)}_\text{increase likelihood of $y_w$} - \underbrace{\nabla_\theta\log\pi(y_l \mid x)}_\text{decrease likelihood of $y_l$}\bigg]\bigg],
\end{multline*}

with $\hat{r}_\theta(x, y) = \beta \log \frac{\pi_\theta(y \mid x)}{\pi_\text{ref}(y \mid x)}$ denoting the reward implicitly determined by the language model $\pi_\theta$ relative to the reference policy $\pi_\text{ref}$ (see Section~\ref{sec:theory} for details). Intuitively, this gradient directs the optimization to raise the likelihood for the favored completions $y_w$, while reducing it for the less-preferred alternatives $y_l$. A central feature is the weighting of examples: the loss gradient emphasizes cases where the implicit reward $\hat{r}_\theta$ assigns a higher value to $y_l$ than to $y_w$, proportionally scaled by $\beta$; this reflects the degree to which the model misorders the completion preferences, moderated by the KL regularization strength. Empirical results highlight the significance of this reweighting mechanism, as omitting the coefficient can result in degenerative behavior of the language model (refer to Appendix Table~\ref{tab:unlikelihood_generations}).

\textbf{DPO overview.} The standard workflow for DPO involves: 1) Drawing completion samples $y_1, y_2 \sim \pi_\text{ref}(\cdot \mid x)$ for each input prompt $x$ and assigning labels based on human judgments, which leads to the construction of an offline preference dataset $\mathcal{D} = \{x^{(i)}, y_w^{(i)}, y_l^{(i)}\}_{i=1}^N$; and 2) training the language model $\pi_\theta$ by minimizing the DPO objective $\mathcal{L}_\text{DPO}$, with respect to the fixed $\pi_\text{ref}$, dataset $\mathcal{D}$, and a selected $\beta$. 
In applied settings, it is often advantageous to leverage existing public preference datasets rather than producing new samples and eliciting human ratings. Preference datasets are typically sampled from $\pi^\text{SFT}$, so whenever possible, we initialize $\pi_\text{ref}$ as $\pi^\text{SFT}$. If $\pi^\text{SFT}$ is absent, we instead set $\pi_\text{ref}$ by performing maximum likelihood estimation over the preferred completions ${(x, y_w)}$, that is, letting ${\pi_\text{ref} = \argmax_{\pi}\mathbb{E}_{x, y_w \sim \mathcal{D}}\left[\log \pi(y_w \mid x)\right]}$. This initialization strategy serves to address the discrepancy between the intractable true reference distribution and the practical $\pi_\text{ref}$ employed by DPO. Additional implementation specifics and hyperparameter settings are described in Appendix~\ref{app:implementation}.

\section{Theoretical Analysis of DPO}
In the present section, we expand on the theoretical underpinnings and interpretation of the DPO approach, clarifying its benefits by comparison to actor-critic techniques such as those employed in RLHF—e.g., PPO~\cite{schulman2017proximal}—and identifying their limitations.

\label{sec:theory}

\subsection{Your Language Model Is Secretly a Reward Model}
DPO achieves policy learning via a single maximum likelihood optimization, without constructing an explicit reward function or carrying out a separate RL process. In fact, the DPO loss (Eq. \ref{eq:main_eq}) coincides with a Bradley-Terry model where rewards are parameterized by $r^*(x, y) = \beta \log\frac{\pi^*_\theta(y \mid x)}{\pi_\text{ref}(y \mid x)}$; consequently, optimizing the parametric model $\pi_\theta$ matches reward model learning as in Eq. \ref{eq:reward_model}, up to a variable substitution. The remainder of this section develops the theoretical rationale for this reparameterization, demonstrates that it does not restrict the family of representable reward models, and establishes that the optimal policy can be recovered exactly. To do so, we introduce a suitable equivalence relation on reward functions.

\begin{definition}
Two reward functions $r(x, y)$ and $r'(x, y)$ are considered equivalent if and only if there exists a function $f$ such that $r(x, y)-r'(x, y) = f(x)$.
\end{definition}

It is straightforward to verify that this definition establishes an equivalence relation, thereby dividing the collection of reward functions into equivalence classes. The following lemmas encapsulate key consequences of this fact:

\begin{lemma}\label{lemma:same_prefrence}
Within the Plackett-Luce framework—including its special case, the Bradley-Terry model—reward functions belonging to the same equivalence class will induce identical preference distributions.
\end{lemma}

\begin{lemma}\label{lemma:same_policy}
    For the constrained RL setting, any two reward functions from a single equivalence class will result in the same optimal policy.
\end{lemma}

The proofs for these statements are elementary and are provided in Appendix \ref{app:lemma1}. The first lemma highlights a classic under-specification property present in models of the Plackett-Luce family~\cite{plackett1975analysis}. Consequently, it is often necessary to incorporate further identifiability constraints to derive well-posed MLE solutions to Eq.~\ref{eq:reward_model}~\cite{bong2022generalized}. According to the second lemma, the optimal policy is invariant across all reward functions from an equivalence class, so it suffices to recover any representative of the optimal equivalence class when optimizing our final objective. The subsequent theorem, proven in Appendix~\ref{app:thm1}, formalizes this insight:

\begin{theorem}\label{thm:main}
    Subject to mild conditions, any reward function that agrees with the Plackett-Luce (and particularly Bradley-Terry) model assumptions can be written as $r(x, y) = \beta \log \frac{\pi(y\mid x)}{\pi_\text{ref}(y\mid x)}$ for some candidate model $\pi(y\mid x)$ and a fixed reference model $\pi_\text{ref}(y \mid x)$.
\end{theorem}

\begin{sproof}
    Let $r(x, y)$ be a reward function with corresponding optimal model $\pi_r(y \mid x)$, as defined in Eq.~\ref{eq:op_policy}. We aim to demonstrate that a representative from $r$'s equivalence class admits the proposed reparameterization. Define the mapping $f$ as
\begin{equation}
    f(r; \pi_\text{ref}, \beta)(x, y) = r(x, y) - \beta\log\sum_{y}\pi_\text{ref}(y\mid x)\exp\left(\frac{1}{\beta}r(x, y)\right)
\end{equation}
Here, $f$ normalizes the reward function by subtracting the log-partition function associated with $\pi_r$, which depends solely on the input $x$. Thus, $f(r; \pi_\text{ref}, \beta)(x, y)$ lies in the equivalence class of $r(x, y)$. Substituting $r$ with the right-hand side of Eq.~\ref{eq:main_eq}, valid for any reward function, yields $f(r; \pi_\text{ref}, \beta)(x, y) = \beta \log \frac{\pi_r(y\mid x)}{\pi_\text{ref}(y\mid x)}$. Consequently, $f$ produces a reward function from $r$'s equivalence class in the desired parameterization, confirming that the proposed reparameterization captures the general case without loss of expressivity.
\end{sproof}

Alternatively, Theorem~\ref{thm:main} can be interpreted as precisely identifying which reward function the DPO reparameterization chooses within each equivalence class—specifically, the function that fulfills:

\begin{equation}\label{eq:lag_p}
     \sum_{y}\underbrace{\pi_\text{ref}(y\mid x)\exp\left(\frac{1}{\beta}r(x, y)\right)}_{=\pi(y\mid x)\text{, using Thm.~\ref{thm:main} reparam.}} = 1,
\end{equation}

This ensures that $\pi(y\mid x)$ constitutes a legitimate probability distribution, meaning all probabilities are nonnegative and collectively sum to one. Referring to Eq.~\ref{eq:op_policy}, we recognize that Eq.~\ref{eq:lag_p} is the partition function associated with the optimal policy defined by the reward $r(x, y)$. The central idea underpinning the DPO algorithm is to impose particular restrictions on the otherwise under-specified Plackett-Luce family (specifically, the Bradley-Terry model), so that while the set of expressible reward models remains unchanged, the optimal policy defined by Eq.~\ref{eq:op_policy} can be computed explicitly for any given prompt $x$.

\subsection{Instability of Actor-Critic Algorithms}

Our analytical framework further enables the diagnosis of instabilities inherent in conventional actor-critic algorithms prevalent in RLHF, such as PPO. In alignment with the RLHF workflow, we concentrate on the RL fine-tuning phase discussed in Section \ref{section:prelims}. Here, connections emerge to the control as inference paradigm \cite{levine2018reinforcement} within the context of the constrained RL problem introduced in \ref{eq:RL}. Consider a parameterized policy model $\pi_{\theta}(y\mid x)$; the typical objective involves minimizing $\mathbb{D}_{\text{KL}}[\pi_{\theta}(y|x) \mid \mid \pi^*(y\mid x)]$, where $\pi^*$—derived from Eq.~\ref{eq:optimum_model}—is the optimal policy corresponding to the reward $r_{\phi}(y, x)$. With appropriate algebraic manipulation, we obtain the following objective:

\begin{equation}\label{eq:AC}
    \max_{\pi_{\theta}}\mathbb{E}_{\pi_{\theta}(y\mid x)}\bigg[\underbrace{r_{\phi}(x, y) -\beta\log\sum_{y}\pi_\text{ref}(y\mid x)\exp\left(\frac{1}{\beta}r_{\phi}(x, y)\right)}_{f(r_{\phi}, \pi_\text{ref}, \beta)} - \underbrace{\beta\log\frac{\pi_{\theta}(y\mid x)}{\pi_\text{ref}(y\mid x)}}_{\text{KL}}\bigg]
\end{equation}

The objective optimized here aligns with that employed in previous studies 
\citep{ziegler2020finetuning, stiennon2022learning, bai2022training, ouyang2022training}, utilizing a reward function $r_{\phi}$ consistent with the DPO formulation. Within this context, the normalization component of $f(r_{\phi}, \pi_\text{ref}, \beta)$ can be interpreted as representing the soft value function for the reference policy $\pi_\text{ref}$. Although this normalization does not alter the optimal policy, omitting it may lead to high variance in policy gradient estimation, potentially destabilizing the learning process. One can substitute the normalization term with an estimated value function, but this approach often encounters optimization challenges. Earlier studies addressed normalization by referencing a human completion as a baseline, effectively providing a single-sample Monte Carlo estimate for the normalization constant. In contrast, the DPO reparameterization constructs a reward function that eliminates the need for explicit baselines.

\section{Experiments}
We conduct a series of empirical studies to assess how well DPO can directly train policies from preference data. Our initial experiments are situated in a controlled text-generation environment and focus on analyzing how effectively DPO balances reward maximization with the minimization of KL-divergence from the reference policy, as compared to prevalent preference optimization algorithms like PPO. Subsequent evaluations consider the application of DPO to larger-scale models and more challenging RLHF domains, such as summarization and conversational tasks. Across these experiments, we observe that DPO generally matches or exceeds the performance of robust baselines—such as RLHF with PPO and choosing the top trajectory out of $N$ samples using a learned reward—requiring little to no hyperparameter adjustment. We precede the presentation of our findings with a description of the experimental design; further implementation specifics are available in Appendix~\ref{app:exp_details}.

\textbf{Tasks.} Our study investigates three open-ended text generation tasks. Across all tasks, methods are trained to optimize a policy based on a preference dataset $\mathcal{D}=\bigl\{x^{(i)}, y_w^{(i)}, y_l^{(i)}\bigr\}_{i=1}^N$. In the \textbf{controlled sentiment generation} setting, $x$ corresponds to the initial segment of a movie review from the IMDb dataset \cite{maas-EtAl:2011:ACL-HLT2011}, and the policy’s goal is to produce $y$ that conveys a positive sentiment. To enable precise assessment, we construct preference pairs of model outputs by utilizing a pre-trained sentiment classifier, ensuring that $p(\text{positive}\mid x,y_w)>p(\text{positive}\mid x,y_l)$. For supervised fine-tuning (SFT), we adapt GPT-2-large by training it to convergence on the training portion of the IMDB review dataset (further details are in App~\ref{app:sentiment_details}). For the \textbf{summarization} task, $x$ is a Reddit forum post, and the objective for the policy is to summarize $y$—extracting the main points. Adopting established approaches, we leverage the Reddit TL;DR summarization dataset \citep{volske-etal-2017-tl} together with human preference annotations obtained by \citeauthor{stiennon2022learning}. Our SFT model, trained on human-curated forum post summaries,\footnote{\url{https://huggingface.co/CarperAI/openai_summarize_tldr_sft}} is fine-tuned with the TRLX \citep{leandro_von_werra_2023_7790115} toolkit for RLHF. The preference data were collected on completions from another SFT model with a comparable training regime. In the case of \textbf{single-turn dialogue}, $x$ represents a human user’s query, which may range from factual inquiries to advice requests. Here, the policy needs to generate an informative and friendly response $y$; we employ the Anthropic Helpful and Harmless dialogue dataset \citep{bai2022training}, featuring 170k interactions between humans and an automated assistant. Each conversation concludes with two completions from a large (unspecified) language model, paired with human annotation indicating the favored response. Since no pre-trained SFT model exists in this scenario, we fine-tune a pre-existing language model exclusively using the preferred responses to obtain our SFT model.

\textbf{Evaluation.} We utilize two complementary evaluation strategies in our experimental framework. To rigorously assess each algorithm's ability to solve the constrained reward maximization task, we first examine their respective frontiers in the space defined by achieved reward versus KL-divergence from the reference policy under a controlled sentiment generation regime. This analysis is feasible since the true reward function—represented here by a sentiment classifier—is accessible. By contrast, in practical applications where the true reward function is unknown, we measure algorithmic performance in terms of \textit{win rate} against a baseline policy. For this, we rely on GPT-4 as a surrogate for human evaluation to judge summary quality in summarization and response helpfulness in the single-turn dialogue setting. In the summarization context, we use ground-truth reference summaries from the test set as baselines; for dialogue, the human-preferred test responses serve as baselines. While recent literature suggests that language models can surpass traditional metrics in automated evaluations \citep{Chen2023ExploringTU}, we supplement this with a human evaluation to validate the reliability of GPT-4 as an evaluator in Sec.~\ref{sec:human-judgments}. The results indicate a strong alignment between GPT-4 assessments and human judgments, with human-to-GPT-4 agreement rates at least matching, and sometimes exceeding, those observed among human annotators themselves.

\textbf{Methods.} Alongside DPO, our comparative study spans a variety of established techniques for training language models according to human preferences. At the simplest end, we examine zero-shot prompting of \textbf{GPT-J} \citep{gpt-j} on summarization, and 2-shot prompting using \textbf{Pythia-2.8B} \citep{biderman2023pythia} on dialogue. The evaluation further encompasses the \textbf{SFT} approach, as well as \textbf{Preferred-FT}, which involves fine-tuning with supervised learning on the preferred completion $y_w$, either generated by the SFT model (for controlled sentiment and summarization) or from a general-purpose language model (for single-turn dialogue). Another related supervised-like approach is \textbf{Unlikelihood}~\citep{welleck2019neural}, where the training objective encourages high probability for $y_w$ while explicitly suppressing $y_l$ through a tunable unlikelihood term modulated by $\alpha\in[0,1]$. Our analysis also considers \textbf{PPO} \citep{schulman2017proximal}, which is optimized with a learned reward model derived from preference data, as well as \textbf{PPO-GT}, a privileged baseline that directly leverages the ground-truth reward function, available only in the controlled sentiment environment. For sentiment generation tasks, we implement two PPO-GT variants: one is an out-of-the-box implementation \cite{leandro_von_werra_2023_7790115}, and the other integrates reward normalization and more extensively tuned hyperparameters to enhance results (these refinements are similarly applied to standard PPO when employing learned rewards). Lastly, we include the \textbf{Best of $N$} baseline, which involves generating $N$ samples from either the SFT model (or Preferred-FT in dialogue) and selecting the highest-ranked completion as scored by the reward model fitted on preference data. Although this method achieves strong results by disentangling reward model accuracy from the PPO optimization process, it remains computationally demanding for even moderate values of $N$, given that it requires generating $N$ candidate outputs per input at inference time.

\subsection{How well can DPO optimize the RLHF objective?}

The standard RLHF framework employs a KL-regularized reward maximization objective, which aims to balance leveraging high reward outcomes with maintaining proximity to the reference policy. Consequently, evaluations of algorithmic performance should jointly consider the magnitude of the obtained reward and the associated KL divergence; it is not sufficient to attain marginally better reward if it results in a substantially larger KL. In Figure~\ref{fig:frontier-tldr-main}, we present the trade-off curve between reward and KL divergence across several algorithms within the sentiment analysis task. For each method, multiple training sessions are conducted, each with distinct settings for policy conservativeness (PPO uses target KL values in $\{3,6,9,12\}$, $\beta$ for DPO set at $\{0.05,0.1,1,5\}$, $\alpha$ for unlikelihood spanning $\{0.05,0.1,0.5,1\}$, and random seeds for preferred-FT), resulting in a total of 22 runs. Following every 100 training iterations up to convergence, we evaluate the policies using a collection of test prompts, measuring both the mean true reward and the mean sequence-level KL divergence\footnote{Summing KL-divergences over all timesteps.} with the reference policy $\text{KL}\left(\pi\mid \mid \pi_\text{ref}\right)$. Our findings indicate that DPO establishes the most effective reward-KL frontier, yielding superior reward while preserving a low KL discrepancy. This outcome is particularly significant for several reasons. Although DPO and PPO target the same objective, DPO demonstrates greater efficiency, yielding a strictly better reward/KL balance than PPO. Furthermore, DPO surpasses PPO in the reward-KL frontier \emph{even under conditions where PPO utilizes true rewards} (PPO-GT).

\subsection{Can DPO scale to real preference datasets?}
\label{sec:dpo-real-datasets}

We proceed to assess the effectiveness of DPO fine-tuning on tasks such as summarization and single-turn dialogue. In summarization, it has been noted that automatic metrics like ROUGE frequently exhibit poor alignment with human judgments~\citep{stiennon2022learning}. Previous research has demonstrated that LMs fine-tuned via PPO using human preference data can produce higher-quality summaries. In our experiments, we compare various approaches by generating completions from the test split of the TL;DR summarization dataset and measuring the average win rate of each method’s outputs against reference completions. Completions for every approach are drawn at different sampling temperatures, ranging from 0.0 to 1.0, and win rates are illustrated in Figure~\ref{fig:frontier-tldr-main} (right). DPO, PPO, and Preferred-FT are all fine-tuned from the same underlying GPT-J SFT model\footnote{\url{https://huggingface.co/CarperAI/openai_summarize_tldr_sft}}. Our results indicate that, at a temperature of 0.0, DPO achieves a win rate of roughly 61\%, surpassing PPO's performance of around 57\% at its best sampling temperature. Furthermore, DPO attains a superior peak win rate when compared with the best of $N$ baseline. It should be highlighted that we did not extensively optimize DPO’s $\beta$ hyperparameter, implying that these outcomes may not fully reflect DPO’s capabilities. Additionally, DPO demonstrates significantly greater stability across different sampling temperatures, in contrast to PPO, whose performance may deteriorate to that of the base GPT-J model as temperature increases. Preferred-FT yields only marginal gains over the SFT baseline. Finally, in Section~\ref{sec:human-judgments}, we conduct direct human preference comparisons between DPO and PPO: human evaluators favored DPO samples at temperature 0.25 in 58\% of instances compared to PPO samples at temperature 0.

For single-turn dialogue evaluation, we analyze the various approaches on a subset of the test portion of the Anthropic HH dataset \citep{bai2022training}, focusing specifically on examples involving one round of human-assistant exchange. In the GPT-4 assessment, the preferred completions from the test set serve as the benchmark for calculating win rates across the methods. Due to the lack of an established SFT baseline, our pipeline begins with a pre-trained Pythia-2.8B, applies Preferred-FT to fine-tune a reference model on the preferred completions (to ensure outputs are model-distributed), and subsequently applies DPO training. Comparisons are also drawn against the best result among 128 Preferred-FT completions—after noting that this "Best of $N$" baseline reaches a performance ceiling at 128 samples for this task (refer to Appendix Figure~\ref{fig:best-of-n})—as well as a 2-shot prompted configuration of the original Pythia-2.8B, with results showing DPO is consistently competitive or superior when optimal temperature settings are used for each method. Additionally, we test an RLHF model, utilizing PPO and trained on the Anthropic HH dataset\footnote{\url{https://huggingface.co/reciprocate/ppo_hh_pythia-6B}}, sourced from a widely recognized implementation\footnote{\url{https://github.com/CarperAI/trlx/tree/main/examples/hh}}; however, no prompt or sampling temperature for this RLHF model outperforms the base Pythia-2.8B. Based on both our TL;DR task findings and the alignment of reward functions, we view Best of 128 as a loose estimator for PPO-level results. Notably, DPO emerges as the sole method capable of enhancing preferred completions on the Anthropic HH dataset with computational efficiency, and matches or exceeds the resource-intensive Best of 128 approach. Moreover, as depicted in Figure~\ref{fig:dialogue-main}, DPO rapidly achieves optimal performance during training.

\subsection{Generalization to a new input distribution}

To further examine how PPO and DPO respond to changes in input distribution, we assess the policies produced by PPO and DPO—trained on Reddit TL;DR summarization—using an out-of-distribution test set: news articles from the CNN/DailyMail dataset \citep{nallapati-etal-2016-abstractive}. We apply the sampling temperatures identified as optimal for TL;DR (0 and 0.25). Table~\ref{tab:ood} summarizes the findings. To measure performance, we determine the win rate of GPT-4 when comparing model outputs to reference summaries, utilizing the same GPT-4 (C) prompt from the Reddit TL;DR evaluation but with the term ``forum post'' replaced by ``news article''. Notably, DPO continues to demonstrate superior performance over PPO by a substantial margin under these new conditions. These results offer preliminary support for the generalizability of DPO policies on par with PPO, despite DPO lacking the supplementary unlabeled Reddit TL;DR prompts available to PPO during training.

\subsection{Validating GPT-4 judgments with human judgments}
\label{sec:human-judgments}
We undertake a human evaluation to test the reliability of GPT-4's output assessments, leveraging results from the TL;DR summarization task and using two distinct prompting strategies. The \textbf{GPT-4 (S)} (simple) prompt requests a judgment on which summary better captures key information from the post. Meanwhile, the \textbf{GPT-4 (C)} (concise) prompt inquires not only about informativeness but also about which summary is more succinct; we include this variant after noting that GPT-4, when following the \textbf{GPT-4 (S)} prompt, tends to favor longer and more repetitive outputs than human raters. The full prompt texts are included in Appendix~\ref{app:prompts}. Our analysis involves three types of comparisons—using the top-performing (DPO, temp. 0.25), lowest-performing (PPO, temp. 1.0), and an intermediate (SFT, temp. 0.25) system—each contrasted against PPO with greedy sampling at its best temperature. This selection ensures a range of summary qualities is represented. Our results indicate that, for both prompts, GPT-4's choices align with those of humans at a rate comparable to inter-human agreement, supporting the validity of GPT-4 as a stand-in for human raters (though for DPO and PPO-1 only, due to constraints on human ratings, do we obtain multiple human judgments). In general, we find the \textbf{GPT-4 (C)} prompt yields win rates that more closely mirror those from human annotators, so we adopt it for the primary analyses in Section~\ref{sec:dpo-real-datasets}. Further information about the human subject study, such as the rater web interface and a roster of participating volunteers, can be found in Appendix~\ref{app:human-study}.

\section{Discussion}
Learning from preferences constitutes an effective and scalable approach for aligning and enhancing the capabilities of language models. In this work, we presented DPO, a straightforward methodology for optimizing language models using preference data, circumventing the need for reinforcement learning. Instead of adapting the preference learning objective to fit into traditional RL frameworks and leveraging conventional RL algorithms, DPO provides a principled correspondence between language model policies and reward structures, facilitating direct optimization to align with human preferences using a cross-entropy loss—thus obviating the need for reinforcement learning and maintaining generality. Remarkably, DPO achieves comparable or superior results relative to established RLHF techniques, such as those leveraging PPO, even with minimal hyperparameter adjustment; as a result, DPO lowers the complexity of training language models from preference data.

\textbf{Limitations \& Future Work.} The findings of our study give rise to several important avenues for further investigation. For instance, the manner in which DPO-trained policies extrapolate to data distributions outside of those seen during training, particularly compared to policies trained from explicit reward signals, warrants deeper analysis. Preliminary evidence indicates that DPO-based policies achieve generalization performance on par with those optimized using PPO, though a broader evaluation is required. Additionally, an open question remains as to whether the use of self-labeling with the DPO policy can similarly exploit prompts without human annotations. Moreover, it is important to explore the manifestations of reward over-optimization within the DPO framework, particularly in light of the modest decline in performance depicted in Figure~\ref{fig:dialogue-main}-right, to determine if it reflects such an effect. While our experiments cover models containing up to 6B parameters, extending DPO to even larger, cutting-edge architectures presents a promising research trajectory. On the evaluation front, our observations reveal that GPT-4's computed win rates are sensitive to the prompt formulation; future research could focus on optimizing prompts to derive higher fidelity assessments from automated evaluators. Ultimately, the principles behind DPO can be extended to a variety of generative modeling domains beyond language modeling from human preferences.